\chapter{TEST PLAN (TP) for Game Rules and Play}

{\centering\selectlanguage{english}\bfseries\color{black}
Version 1.1, April 2014
\par}

\ 

\bigskip

{\centering\bfseries\color{black}
TEST PLAN (TP)}

{\centering\selectlanguage{english}\bfseries\color{black}
FOR
\par}


\bigskip

{\centering\selectlanguage{english}\bfseries\color{black}
Swords and Sorcery
\par}


\bigskip


\bigskip


\bigskip

%\begin{figure}
%\centering
%\includegraphics[width=3.4354in,height=0.6126in]{uislogan.png}
%\end{figure}

\bigskip


\bigskip


\bigskip


\bigskip

{\centering\selectlanguage{english}\bfseries\color{black}
Version 1
\par}

{\centering\selectlanguage{english}\bfseries\color{black}
April 30, 2014
\par}


\bigskip


\bigskip

{\centering\selectlanguage{english}\bfseries\color{black}
Prepared for:
\par}

{\centering\selectlanguage{english}\bfseries\color{black}
Clinton Jeffery
\par}


\bigskip


\bigskip

{\centering\selectlanguage{english}\bfseries\color{black}
Prepared by:
\par}

{\centering\selectlanguage{english}\bfseries\color{black}
Tao Zhang, Cameron Simon, Matthrew Brown, \\Tyler Jaszkowiak, Ian Westrope, ChiHsiang Wang \\(Rule Team)
\par}

{\centering\selectlanguage{english}\bfseries\color{black}
University of Idaho
\par}

{\centering\selectlanguage{english}\bfseries\color{black}
Moscow, ID \ 83844-1010
\par}

{\centering\selectlanguage{english}\bfseries\color{black}
CS383 TPD
\par}

\pagebreak

{\centering\selectlanguage{english}\bfseries\color{black}
RECORD OF CHANGES (Change History)
\par}

\begin{flushleft}
\tablehead{}
\begin{supertabular}{|m{0.5462598in}|m{0.6712598in}|m{1.4212599in}|m{0.23375985in}|m{1.7962599in}|m{0.7337598in}|m{0.6295598in}|}
\hline
~

\centering {\selectlanguage{english}\bfseries\color{black} Change}\par

\centering {\selectlanguage{english}\bfseries\color{black} Number}\par

~
 &
~

\centering \selectlanguage{english}\bfseries\color{black} Date completed
&
~

\centering {\selectlanguage{english}\bfseries\color{black} Location of
change }\par

\centering \selectlanguage{english}\bfseries\color{black} (e.g., page or
figure \#) &
~

\centering {\selectlanguage{english}\bfseries\color{black} A}\par

\centering \selectlanguage{english}\bfseries\color{black} M\newline
D  &
~

\centering {\selectlanguage{english}\bfseries\color{black} Brief
description }\par

\centering \selectlanguage{english}\bfseries\color{black} of change &
~

\centering \selectlanguage{english}\bfseries\color{black} Approved by
(initials) &
~

\centering {\bfseries\color{black} Date }\par

\centering\arraybslash\bfseries\color{black}
approved\\

 &

 &

 &

 &

 &

 &

\\\hline
1
 &
4/30/2014
 &
Document
 &
A
 &
Original Document Created
 &
MB
 &
4/30/2014
\\\hline
~
 &
~
 &
~
 &
~
 &
~
 &
~
 &
~
\\\hline
~
 &
~
 &
~
 &
~
 &
~
 &
~
 &
~
\\\hline
~
 &
~
 &
~
 &
~
 &
~
 &
~
 &
~
\\\hline
~
 &
~
 &
~
 &
~
 &
~
 &
~
 &
~
\\\hline
~
 &
~
 &
~
 &
~
 &
~
 &
~
 &
~
\\\hline
~
 &
~
 &
~
 &
~
 &
~
 &
~
 &
~
\\\hline
~
 &
~
 &
~
 &
~
 &
~
 &
~
 &
~
\\\hline
~
 &
~
 &
~
 &
~
 &
~
 &
~
 &
~
\\\hline
~
 &
~
 &
~
 &
~
 &
~
 &
~
 &
~
\\\hline
~
 &
~
 &
~
 &
~
 &
~
 &
~
 &
~
\\\hline
~
 &
~
 &
~
 &
~
 &
~
 &
~
 &
~
\\\hline
~
 &
~
 &
~
 &
~
 &
~
 &
~
 &
~
\\\hline
~
 &
~
 &
~
 &
~
 &
~
 &
~
 &
~
\\\hline
~
 &
~
 &
~
 &
~
 &
~
 &
~
 &
~
\\\hline
~
 &
~
 &
~
 &
~
 &
~
 &
~
 &
~
\\\hline
~
 &
~
 &
~
 &
~
 &
~
 &
~
 &
~
\\\hline
~
 &
~
 &
~
 &
~
 &
~
 &
~
 &
~
\\\hline
~
 &
~
 &
~
 &
~
 &
~
 &
~
 &
~
\\\hline
~
 &
~
 &
~
 &
~
 &
~
 &
~
 &
~
\\\hline
~
 &
~
 &
~
 &
~
 &
~
 &
~
 &
~
\\\hline
\end{supertabular}
\end{flushleft}
{\selectlanguage{english}\color{black}
A - ADDED \ M - MODIFIED \ D -- DELETED}

{\centering\selectlanguage{english}\bfseries\color{black}
Swords and Sorcery
\par}

\pagebreak

{\centering\selectlanguage{english}\bfseries\color{black}
TABLE OF CONTENTS
\par}

{\selectlanguage{english}\bfseries\color{black}
Section\ \ Page}

\setcounter{tocdepth}{9}
\renewcommand\contentsname{}
\tableofcontents

\setcounter{page}{1}\pagestyle{Convertiv}

\section[IDENTIFIER]{\selectlanguage{english}\bfseries\color{black} TEST PLAN IDENTIFIER}

{\selectlanguage{english}\color{black}
This test plan has no unique identifier other than Spring 2014 CS 383 sworsorc. This particular test plan is in conjunction to the rules and game play group test plans. There are many test plans for this group due to the size and span of the group's contents. This test plan, for example, will be known as the Units test plan.}

\section[REFERENCES]{\selectlanguage{english}\bfseries\color{black}
REFERENCES}
%{\selectlanguage{english}\itshape\color{black}

{\selectlanguage{english}\color{black}
Use Cases and State diagrams are available at https://github.com/cjeffery/sworsorc/tree/master/doc }


\section[INTRODUCTION]{\bfseries\color{black} INTRODUCTION}

{\selectlanguage{english}\color{black}
The purpose of this test plan is to state the processes used by Game Rules and Play Team in testing the Spells and Characters for the Swords \& Sorcery project.(Tao, Cameron)
\newline
\newline
This test plan covers the creation and implementation of the Moveableunit class and its sub class Armyunits. (Matt)
\newline
\newline
This test plan is also to provide as much coverage as possible to the methods and data of the Scenario class by verifying the results of execution against expectations through a series of automated unit tests and a manual GUI test. (Tyler)
}

\section[TEST ITEMS]{\bfseries\color{black} TEST ITEMS}

{\selectlanguage{english}\color{black}

The army units will be tested for the proper member variable values as well as proper results from the member function of the Moveableunit and Armyunit classes. Variables will tested at creation and members variable will be tested for proper results when applicable. Some example tests will be that the location member variable is properly set during movement. Also another type of test that will be preformed is that conjured units are properly created and and placed in the proper place on the game board. 
}

{\selectlanguage{english}\color{black}
\subsection{Character}
\begin{itemize}
\item Class Interfaces
\item Class Interactions
\item Spells Implementation
\item Character 
\end{itemize}

\subsection{Scenario}
These are things you intend to test within the scope of this test
plan. Essentially, something you will test, a list of what is to be
tested. This can be developed from the software application
inventories as well as other sources of documentation and information.
\newline
\newline
The data read by the Scenario test will be read from one of the simple 
scenario configuration files and then verified against expectations through 
the class's accessor methods. These data items include 
\newline
\begin{itemize}
\item The scenario's name
\item Number of players
\item Game length
\item The blue sun's initial position
\item The names of armies in the scenario
\item The controlling players of these armies
\item The setup order of these armies
\item The movement order of these armies
\item Names of nations within these armies
\item Names of the neutrals in the scenario
\item The provinces controlled by a nation
\item The characters within a nation
\item The units within a player nation
\item The units within a neutral nation
\item The races of both neutrals and player nations
\item The reinforcement and replacement description strings
\item Data related to where a neutral is leaning toward
\item Whether or not a neutral accepts human sacrifice
\end{itemize}

Unfortunately, the most complex functionality of the Scenario class 
cannot be tested by automated unit tests. The unit pool populator requires 
a manual check.
}

\section[SOFTWARE RISK ISSUES]{\bfseries\color{black} SOFTWARE RISK ISSUES}
{\selectlanguage{english}\color{black}

{\selectlanguage{english}\color{black}
\begin{itemize}
\item Improper casting of units
\item Unit ID's interpreted incorrectly
\item Depends on Java's JSON reader and the programmer's understanding of it
\item Complex data structures such as a map of maps
\item Complex loops to iterate over data structures
\item Poor documentation surrounding some of these iterators
\end{itemize}
}

\section[FEATURES TO BE TESTED]{\bfseries\color{black} FEATURES TO BE TESTED}

{\selectlanguage{english}\color{black}
\subsection{Scenario}
From the user's perspective, much of this data reading occurs under the hood. 
In all cases except for populating the unit pool, the Scenario class does not 
manipulate any pieces of the rest of the project. Other components read the 
data from the Scenario class. Therefore, while the solar configuration relies 
on a properly-working Scenario class, this is not apparent to the user because 
solar configuration is also handled by SolarConfig and HUDInitializer classes.
\newline
\newline
Therefore, the only visible component being tested by this plan is the placement 
of units and characters into the correct provinces of the map.

\subsection{Movement}

\begin{itemize}
\item conjured unit appear when casted
\item units are properly represented on HUD
\item movement location is correct on map
\item moral status is properly updated after combat
\end{itemize}

\subsection{Character, Spells}
\begin{itemize}
\item Select character on GUI
\item Select Spell
\item Effects of Partial Spells
\item Manna Costing
\end{itemize}
}

\section[FEATURES NOT TO BE TESTED]{\bfseries\color{black} FEATURES NOT TO BE TESTED}

{\selectlanguage{english}\color{black}
This is a listing of what is NOT to be tested from both the Users
viewpoint of what the system does and a configuration
management/version control view. This is not a technical description
of the software, but a USERS view of the functions.
\newline
\newline
All data is verified, but testing the Scenario class alone cannot ensure 
that it reaches the HUD successfully. Integration tests are required outside 
of this plan for information such as solar configuration, move order, setup order, 
game length, and diplomacy.
\newline
\newline
What is not tested is the count of the units on the map that correspond to the scenario loaded. Also the type of unit is not test, this is assumed correct. 
}

\section[APPROACH]{\bfseries\color{black} APPROACH}
{\selectlanguage{english}\color{black}

{\selectlanguage{english}\color{black}

We will use JUnit tests as well as manual tests to determine the effectiveness and accuracy of the methods and variables tested. 

}

\section[ITEM PASS/FAIL CRITERIA]{\bfseries\color{black}
	 ITEM PASS/FAIL CRITERIA}
{\selectlanguage{english}\color{black}
Items will receive a PASS once the variables/methods have been test multiple times and under different circumstances and the results of the test were within expected parameters. A FAIL will be received if the test results are not within expected parameters. 
}

\subsection{Unit Testing Pass/Fail Criteria}
\subsection{Intergration Testing Pass/Fail Criteria}
\subsubsection{Scenario}
The automated tests must all pass and result in 100% instruction coverage for 
each of the class's accessor methods. One of the nations in the scenario must 
have the correct number and types of units loaded into the provinces it controls.
\subsubsection{Character}
\begin{center}
\begin{supertabular}{|p{4.5cm}|p{4.5cm}|p{4.5cm}|}
\hline
    \textbf{Rule Description}
    &
    \textbf{Test Description}
    &
    \textbf{Expected Result}
\\\hline
    TBD
    &
    createCharacter function is called with 
    specific character name from CharacterMaker.java.
    &
    A new character object is created and returned 
    that contains all of the specified characters information. 
\\\hline    
\end{supertabular}
\end{center}

\subsubsection{Spells}
\begin{center}
\begin{supertabular}{|p{4.5cm}|p{4.5cm}|p{4.5cm}|}
\hline
    \textbf{Rule Description}
    &
    \textbf{Test Description}
    &
    \textbf{Expected Result}
\\\hline
    Character with Power Level should have a spell book.
    &
    Generate a spell book for the character.
    &
    A frame with a list of spells should be shown on the screen.
\\\hline 
    Show spell description.
    &
    Click on the Spell button.
    &
    A frame will be displayed with all information about the spell.
\\\hline   
    A target need to be selected for most of the spells.
    &
    Click on cast button on the frame of displaying information about the spell. 
    Then select a target by right click the target on the map.
    &
    A target is selected to cast the spell.
\\\hline
\end{supertabular}
\end{center}





\section[SUSPENSION CRITERIA]{\bfseries\color{black}
	 SUSPENSION CRITERIA AND RESUMPTION REQUIREMENTS}
{\selectlanguage{english}\color{black}

{\selectlanguage{english}\color{black}
Tests will be suspended if a bigger piece of code that directly or indirectly effects the code is compromised. Tests will also be suspended once a FAIL has be received for that test. The tests will be resumed once the effected code fixed. 
}

\section[TEST DELIVERABLES]{\bfseries\color{black} TEST DELIVERABLES}

{\selectlanguage{english}\color{black}
\begin{itemize}
\item Test plan document
\item Test cases
\item Relevant error logs or problem reports
\item JUnit test cases
\item A report of passes/failures from JUnit
\item A coverage report from JaCoCo
\end{itemize}
}

\section[REMAINING TEST TASKS]{\bfseries\color{black} REMAINING TEST TASKS}

{\selectlanguage{english}\color{black}
This is a multi-group project in which many test plans are created for the respective group. Also since the full project as not been completed there will be more tests developed in the future. 
\newline
\newline
All other classes in the project require some level of testing as well. 
In particular, the HUD and SolarConfig related classes directly affect 
the performance of Scenario.
}

\section[ENVIRONMENTAL NEEDS]{\bfseries\color{black} ENVIRONMENTAL NEEDS}

{\selectlanguage{english}\color{black}
There are no special requirements for this test plan. 
}

\section[STAFFING AND TRAINING NEEDS]{\bfseries\color{black} STAFFING AND TRAINING NEEDS}
{\selectlanguage{english}\color{black}

{\selectlanguage{english}\color{black}
Training in JUnit test is required to understand sections of the code. 
\newline
\newline
Understading of NetBeans, Java, JUnit, and JaCoCo.
}

\section[RESPONSIBILITIES]{\bfseries\color{black} RESPONSIBILITIES}
{\selectlanguage{english}\color{black}

{\selectlanguage{english}\color{black}
This is a co-operative, no one person is in charge. Tests are divided by group type developers. 
}

\section[SCHEDULE]{\bfseries\color{black} SCHEDULE}
{\selectlanguage{english}\color{black}


{\selectlanguage{english}\color{black}
Realistically most tests will not be completed due to the rapid approach of the end of the semester.  
}

\section[PLANNING RISKS AND CONTINGENCIES]{\bfseries\color{black} PLANNING RISKS AND CONTINGENCIES}

{\selectlanguage{english}\color{black}
Too many tests will cause a halt in development and since the project wont be completed by the end of the semester might as well get as much working functionality as possible. 
}

\section[APPROVALS]{\bfseries\color{black} APPROVALS}
{\selectlanguage{english}\color{black}
Dr. J is the evil genius behind the project and is the only one capable of approving this process.  (Matt)
\newline
\newline
Dr. Jeffery has the final word. However, I imagine that the team 
generating the test documentation will decide which tests can be 
approved and whether to continue.(Tyler)
}

\section[GLOSSARY]{\bfseries\color{black} GLOSSARY}
{\selectlanguage{english}\color{black}
JaCoCo: The Java Code Coverage plugin for NetBeans}



\clearpage\setcounter{page}{1}\pagestyle{Convertviii}
\section[APPENDIX A. \ [insert name
here{]}]{\selectlanguage{english}\bfseries\color{black} APPENDIX A.
\ [insert name here]}
{\selectlanguage{english}\itshape\color{black}
Include copies of test examples, etc. supplied or
derived from the customer. \ Appendices are labeled A, B, {\dots}n.
\ \ Reference each appendix as appropriate in the text of the document.
}

{\selectlanguage{english}\color{black}
\ [ insert appendix A here ]}

\clearpage\setcounter{page}{1}\pagestyle{Convertix}
\section[APPENDIX B. \ [insert name
here{]}]{\selectlanguage{english}\bfseries\color{black} APPENDIX B.
\ [insert name here]}

\bigskip

{\selectlanguage{english}\color{black}
[ insert appendix B here ]}


\bigskip
